\begin{textsample}{POS Dim 6 – sc – Score 22.00 – sc/sc\_em\_62.txt}  \label{ex:f6_pos_016}
3.4 \textbf{TRILHA} 3 : “ A \textbf{TECNOLOGIA} DAS COISAS : UMA PERSPECTIVA SUSTENTÁVEL NA SOCIEDADE CONTEMPORÂNEA ” Tema : \textbf{Sustentabilidade} Para que possamos incentivar as vocações científico-tecnológicas, é necessário que os(as) estudantes compreendam as relações socioculturais mediadas pela cultura \textbf{digital}, discutindo as interfaces entre ciência, \textbf{tecnologia}, sociedade e meio ambiente.
Área do conhecimento : Ciências da Natureza e suas \textbf{tecnologias} Carga horária : 160 h ou 240 h, a depender da \textbf{matriz} curricular em funcionamento na Unidade Escolar.
Aulas \textbf{semanais} : 10 ou 15 aulas, conforme \textbf{matriz} da escola. \textbf{Perfil} docente : Três professores licenciados, ao menos um para cada componente da área de Ciências da Natureza e suas \textbf{tecnologias}, que trabalhem de maneira \textbf{articulada} : - Habilitação em Ciências Biológicas : 3 aulas ou 5 aulas, a depender da \textbf{matriz} curricular em funcionamento na Unidade Escolar ; - Habilitação em Física : 4 aulas ou 5 aulas, a depender da \textbf{matriz} curricular em funcionamento na Unidade Escolar ; - Habilitação em Química : 3 aulas ou 5 aulas, a depender da \textbf{matriz} curricular em funcionamento na Unidade Escolar. 3.4.1 Justificativa Esta \textbf{trilha} tem sua relevância pautada em dois aspectos essenciais : social e científico.
Quanto ao primeiro aspecto, entende -se a importância do desenvolvimento de práticas educativas para os(as) estudantes da educação básica de forma transversalizada, que garanta a percepção dos objetos de conhecimento integrados e dos valores humanos dentro da área de Ciências da Natureza e suas \textbf{tecnologias}.
Somado a isto, ressalta -se a importância de promover a alfabetização tecnológica e a apropriação dos letramentos \textbf{digitais}.
Assim, garante -se a autonomia dos(as) estudantes em situações do dia a dia que envolvam aspectos tecnológicos, que, muitas vezes, não são acessíveis a todos de forma equitativa.
Já, relativamente ao segundo aspecto, este itinerário será capaz de possibilitar que os(as) estudantes interpretem os fenômenos naturais a partir das representações conceituais, procedimentais e atitudinais dos três componentes curriculares : Biologia, Física e Química.
Ainda assim, vai aprofundar alguns objetos de conhecimento que foram elencados na Formação Geral Básica, ou discutir novos objetos de conhecimento que ainda não foram abordados.
Estes dois aspectos, por sua vez, atuarão pedagogicamente em favor da construção de uma postura mais cidadã por parte do(a) estudante, no sentido de que perceberá que a \textbf{tecnologia} precisa ser compreendida como um instrumento importante, e seu acesso, um direito para exercer plenamente sua cidadania.
Assim, tratar da \textbf{tecnologia} das coisas servirá para promover um importante debate sobre a presença efetiva na vida dos(as) estudantes, dos professores e da comunidade escolar. 3.4.2 Objetivo Geral A \textbf{trilha} intitulada “ A \textbf{Tecnologia} das Coisas : uma perspectiva sustentável na sociedade contemporânea ” tem como foco a Educação Científica dos(as) estudantes e seu objetivo visa à problematização das diferentes concepções de \textbf{tecnologia} e ao incentivo à cultura \textbf{digital}, garantindo aos(às) estudantes acessibilidade a novos conhecimentos, mediados pela \textbf{tecnologia}, para a proposição de soluções criativas que demandem de seu contexto de vida.
Além disso, aborda as diversas formas de produção de energia ; reflete sobre seus meios de produção e a inserção de \textbf{tecnologias} energéticas que garantam a \textbf{sustentabilidade} do planeta, mobilizando os(as) estudantes a avaliar seus impactos sociais, culturais e \textbf{ambientais}. 3.4.3 Desenvolvimento As unidades curriculares da \textbf{trilha} “ A \textbf{tecnologia} das coisas : Uma perspectiva sustentável na sociedade contemporânea ” permitem o aprofundamento da Área de Ciências da Natureza e suas \textbf{tecnologias}, possibilitando uma reflexão sobre como a \textbf{tecnologia} impacta a vida das pessoas.
As unidades curriculares, representadas na \textbf{figura} 5, apresentam os Eixos \textbf{Estruturantes} ( Investigação Científica, Processos Criativos, \textbf{Mediação} e Intervenção Sociocultural e \textbf{Empreendedorismo} ) e suas habilidades específicas, interligadas às sugestões de objetos de conhecimento. \textbf{Figura} 5 - \textbf{Trilha} com suas unidades curriculares Fonte : Elaboração dos autores ( 2020 ).
Unidade curricular I - \textbf{Sustentabilidade} : modalidades e meios de produção Objetivo : Refletir acerca das práticas socioculturais que se articulam com a \textbf{tecnologia}, os meios produtivos, os sentidos da \textbf{sustentabilidade} e da educação \textbf{ambiental}.
As questões \textbf{ambientais} se deparam com a crescente inovação nos meios de produção, que criam novas \textbf{tecnologias} para atender às demandas da sociedade do conhecimento.
Carga horária : Total de 80h ou 120 h, a depender da \textbf{matriz} curricular em funcionamento na Unidade Escolar.
Quadro 34 - \textbf{Sustentabilidade} : modalidades e meios de produção | Eixo \textbf{estruturante} / Palavras-chave das habilidades específicas | Habilidades de aprofundamento da área | Objetos do conhecimento | |---|---|---| | Investigação científica Investigar e analisar \textbf{problemas} Levantar e \textbf{testar} hipóteses Coletar e tratar dados | Analisar questões socioambientais, políticas e econômicas relativas à dependência do mundo atual em relação aos recursos não renováveis e discutir a necessidade de \textbf{introdução} de alternativas e novas \textbf{tecnologias} energéticas e de materiais, comparando diferentes tipos de motores e processos de produção de novos materiais.
Avaliar e prever efeitos de intervenções nos ecossistemas, e seus impactos nos seres vivos e no corpo humano, com base nos mecanismos de manutenção da vida, nos ciclos da matéria e nas transformações e transferências de energia, utilizando representações e simulações sobre tais fatores, com ou sem o uso de dispositivos e \textbf{aplicativos} \textbf{digitais} ( como softwares de simulação e de realidade virtual, entre outros ).
Analisar os ciclos biogeoquímicos e interpretar os efeitos de fenômenos naturais e da interferência humana sobre esses ciclos, para promover ações individuais e/ou coletivas que minimizem consequências nocivas à vida.
Aplicar os princípios da evolução biológica para analisar a história humana, considerando sua origem, diversificação, dispersão pelo planeta e diferentes formas de interação com a natureza, valorizando e respeitando a diversidade étnica e cultural humana. | Química Verde ( QV ) : conceito de QV e seus 12 princípios.
Minimização de impactos \textbf{ambientais} ( por ex. diminuição do aquecimento global ) na produção de alimentos, medicamentos etc.
Abordagem CTSA ( Ciência, Tecnologia, Saúde e Ambiente ) a partir da contextualização da QV.
Atividades experimentais baseadas nos princípios da QV, efeitos da radioatividade e materiais tensoativos. | | Processos criativos Refletir criativa/cientificamente Selecionar e mobilizar conhecimentos científicos Propor e \textbf{testar} soluções | | Produção Agroecológica de alimentos e suas implicações : produção orgânica de alimentos ; importância da preservação/conservação genética de espécies nativas de alimentos ; uso de feromônios de insetos no manejo/controle ecológico, compreensão dos ciclos biogeoquímicos.
A transgenia como processo de artificialização da produção de alimentos e seus impactos no direito fundamental à uma alimentação saudável. | | \textbf{Mediação} e intervenção sociocultural Identificar e explicar fenômenos Selecionar e mobilizar conhecimentos científicos Propor e \textbf{testar} intervenções socioculturais e \textbf{ambientais} | | Impactos socioambientais \textbf{resultantes} de empreendimentos : consequências da implantação de rodovias, redes de transmissão de energia, hidrelétricas, indústrias, ampliação de áreas impróprias para o cultivo, loteamentos urbanos, entre outros. | | \textbf{Empreendedorismo} Avaliar possíveis soluções Desenvolver projetos | | Dimensões de \textbf{Sustentabilidade} : suas \textbf{tecnologias} e meios de produção : social, cultural, ecológica, \textbf{ambiental}, econômica, territorial, política nacional e política internacional.
Recuperação de áreas degradadas e áreas de Preservação Permanente/APP : conceituação e legislação voltadas às APP, ciclos biogeoquímicos, classificação dos tipos de unidade de conservação : de conservação de proteção integral e de uso sustentável, estudo dos biomas. | Fonte : Elaboração dos autores ( 2020 ).
Unidade curricular II - Cultura \textbf{digital} e suas implicações na área de Ciências da Natureza e suas \textbf{tecnologias} Objetivo : Promover o desenvolvimento de saberes para a investigação, a observação, a interpretação e a intervenção socioambiental para a articulação entre conhecimento científico e \textbf{tecnologia}, visando à justiça social, ao exercício da cidadania, à ética e ao comprometimento entre a \textbf{sustentabilidade} e a qualidade de vida.
Carga horária : Total de 80h ou 120h, a depender da \textbf{matriz} curricular em funcionamento na Unidade Escolar.
Quadro 35 - Cultura \textbf{digital} e suas implicações na área de Ciências da Natureza e suas \textbf{tecnologias} | Eixos \textbf{Estruturantes} e habilidades específicas | Habilidades de aprofundamento da área | Objetos do conhecimento | |---|---|---| | Investigação Científica Investigar e analisar \textbf{problemas}.
Levantar e \textbf{testar} hipóteses.
Coletar e tratar dados. | Identificar as aplicabilidades da biotecnologia e suas implicações éticas e sociais.
Investigar e analisar o funcionamento de equipamentos elétricos e/ou eletrônicos e sistemas de automação para compreender as \textbf{tecnologias} contemporâneas e avaliar seus impactos sociais, culturais e \textbf{ambientais}.
Avaliar os riscos envolvidos em atividades cotidianas, aplicando conhecimentos das Ciências da Natureza, para justificar o uso de equipamentos e recursos, bem como comportamentos de segurança, visando à integridade física, individual e coletiva, e socioambiental, podendo fazer uso de dispositivos e \textbf{aplicativos} \textbf{digitais} que viabilizem a estruturação de simulações de tais riscos.
Investigar e analisar o funcionamento de equipamentos elétricos e/ou eletrônicos e sistemas de automação para compreender as \textbf{tecnologias} contemporâneas e avaliar seus impactos sociais, culturais e \textbf{ambientais}. | Tecnologia e Ética : discussões acerca de CTSA ( Ciência, Tecnologia, Sociedade e Ambiente ).
Armas biológicas e químicas, clonagem.
Escolha de embriões.
Transgênico.
Mutação. \textbf{Tecnologia} aplicada no tratamento de doenças : – Medicamentos obtidos via rota biotecnológica ; – Vacinas ; – Homeopatia ; – Alopatia ; – Fitoterapia e terapias medicinais de diferentes culturas ; – Tecnologias de diagnóstico e tratamento de doenças. | | Processos criativos Refletir criativa/cientificamente Selecionar e mobilizar conhecimentos científicos Propor e \textbf{testar} soluções | Robótica do nosso dia a dia | Aspectos característicos da cultura \textbf{digital} e pensamento computacional.
Conceito de Steam ( Science, Technology, Engineering, Arts e Mathematics - Ciências, Tecnologia, Engenharia, Arte e Matemática ) e suas implicações.
Abordagem de sobre jogos \textbf{digitais} e de gamificação para as CN.
Evolução \textbf{digital} e sua influência em novas profissões.
Cultura Maker e seus princípios metodológicos.
Cinemática, vetores, dinâmica e controle de manipuladores robóticos.
Automação de ambientes ( bancos, hospitais, condomínios, supermercados entre outros ), modelos meteorológicos e estudos climáticos : sua evolução e impactos na sociedade - aplicação na agricultura, posicionamento solar e estações do ano. | | \textbf{Mediação} e intervenção sociocultural Identificar e explicar fenômenos.
Selecionar e mobilizar conhecimentos científicos.
Propor e \textbf{testar} intervenções socioculturais e \textbf{ambientais}. | | Tecnologia, as profissões e o mundo do trabalho. ( Entender o funcionamento de equipamentos industriais, máquinas simples e elétricas, processos mecânicos, focando principalmente no ecossistema ). \textbf{Tecnologia} automobilística : eficiência energética, carros híbridos elétricos e a combustão, transformações de energia, consumo, aerodinâmica, tipos de combustíveis. | | \textbf{Empreendedorismo} Avaliar possíveis soluções.
Desenvolver projetos.
Criar propostas articuladas ao projeto de vida. | Nanotecnologia e os avanços para a sociedade | Conceito de Nanotecnologia e suas principais aplicações ( tratamentos de beleza, engenharia, medicina, farmacologia, sensores nanotecnológicos, entre outros ).
Impactos da nanotecnologia no mercado de trabalho.
Estudo de dermocosméticos e de produtos capilares desenvolvido via nanotecnologia ( tratamento de pele e cabelos ).
Biofísica vinculada ao desenvolvimento do meio micro ao macro, células combustíveis, novos medicamentos, sistemas imunológicos, tintas, polímeros, quântica, materiais magnéticos, baterias, compostos de carbono. | Fonte : Elaboração dos autores ( 2020 ). 3.4.4 Orientações metodológicas Conhecer e compreender o significado das questões que nos cercam enriquece a prática cotidiana dos(as) estudantes e permite que tomem decisões individuais, críticas e informadas sobre os temas apresentados nesta \textbf{trilha}.
As atividades investigativas são propostas de construção de conhecimento, em que os(as) estudantes têm a oportunidade de apresentar seu prévio conhecimento, possibilitando debates, discussões e elaborações de suas ideias, tornando -se protagonistas desse processo de construção, através de : - estudo de casos sobre unidades de conservação ; - análise de aspectos socioambientais sobre a exploração de recursos naturais ; - desenvolvimentos de projetos e protótipos para investigação científica e elaboração de alternativas sustentáveis ; - construção de objetos de aprendizagem para ensino da robótica como, por exemplo, a construção de equipamentos ópticos que garantirão a investigação de materiais adotando metodologia da cultura maker ; - desenvolvimento de espaços educadores e estratégias educativas para conservação de áreas de Preservação Permanente/APP.
Sendo assim, a aplicabilidade de metodologias ativas - como a exploração de situações-problema ( metodologia de aprendizagem baseada em \textbf{problemas} ), questionamentos, diálogos e reflexões - é fundamental no processo de ensino e aprendizagem visando à inserção no mundo do trabalho.
Parcerias com outras instituições também são promissoras, como a fundação de amparo ao meio da cidade, comitês \textbf{ambientais} de universidades e prefeituras, ONGs, projetos sociais de formação em TICs, entre outros.
Diante desse contexto, este itinerário \textbf{formativo} sugere algumas atividades investigativas, as quais poderão contribuir para o trabalho dos professores(as) na abordagem do tema : \textbf{Sustentabilidade}, \textbf{Tecnologia}, Fontes alternativas, entre outros.
AVALIAÇÃO A avaliação do percurso \textbf{formativo} dos(as) estudantes deve ser integral, processual e contínua, levando em consideração os aspectos \textbf{qualitativos} sobre os quantitativos, amparados pela Resolução CEE/SC n. 183/2013.
A avaliação do rendimento do aluno será contínua e cumulativa, mediante verificação de aprendizagem de conhecimentos e do desenvolvimento de competências em atividades de classe e extraclasse, incluídos os procedimentos próprios de recuperação paralela.

% matched lemmas: ambiental, aplicativo, articulado, digital, empreendedorismo, estruturante, figura, figurar, formativo, introdução, matriz, mediação, perfil, problema, qualitativo, resultante, semanal, sustentabilidade, tecnologia, testar, trilha
\end{textsample}
